\section{Matching moments and a stable old filtration}\label{sec:moments}

For \(m,N\ge1\), the functional unary system
\(\mathcal F^{\mathrm{fun}}_{m,N}\) has variables \(X_{ij}\) and generators
\begin{align*}
 &X_{ij}^2-X_{ij},\\
 &X_{ij}X_{ij'} &&(j\ne j'),\\
 &X_{ij}X_{i'j} &&(i\ne i'),\\
 &\rho_i-1,\qquad \rho_i=\sum_{j=1}^N X_{ij}.&
\end{align*}
The row exclusions in the second line are part of this auxiliary system.
They are not inferred from a weaker unary encoding.

A matching \(T\) is a set of cells with distinct rows and distinct columns;
write \(X_T=\prod_{(i,j)\in T}X_{ij}\).
The empty matching corresponds to the constant one.

\begin{lemma}[Complete moment equations]\label{lem:marginals}
A linear functional on \(\Pol_{\le B}\) annihilating
\(\Ideal_B(\mathcal F^{\mathrm{fun}}_{m,N})\) is exactly a choice of
matching moments \(z_T\), \(|T|\le B\), satisfying
\begin{equation}\label{eq:marginals}
 \sum_{j\notin\cols(T)}z_{T\cup\{(i,j)\}}=z_T
 \quad\text{for }|T|<B,\ i\notin\rows(T).
\end{equation}
Its value on an ordinary monomial is the moment of its squarefree support
if that support is a matching, and is zero otherwise. Normalization is the
additional condition \(z_\varnothing=1\).
\end{lemma}
\begin{proof}
Boolean equations reduce powers within the original degree.
A squarefree monomial whose support has a repeated row or column contains
one of the exclusion generators as a factor. Thus every polynomial reduces,
within its degree, to a linear combination of matching monomials.
Conversely, the stated rule for evaluating ordinary monomials annihilates
every legal Boolean or exclusion multiple.

It remains to check row equations. In a row-equation multiple
\(q(\rho_i-1)\) of degree at most \(B\), we may first put \(q\) into the
matching normal form: any error is a Boolean or exclusion consequence
through \(B-1\), and multiplication by \(\rho_i-1\) keeps its certificate
within \(B\). For a matching multiplier \(X_T\), if row \(i\) is already
occupied, the product reduces to zero. Otherwise its reduction is
\[
 \sum_{j\notin\cols(T)}X_{T\cup\{(i,j)\}}-X_T.
 \]
These are precisely \eqref{eq:marginals}, and they exhaust the constraints.
No positivity or multiplicativity is imposed on the functional.
\end{proof}

Let \(\Delta_{s,N}\) be the chessboard complex whose faces are matchings on
an \(s\)-row, \(N\)-column board. The following consequence of the
chessboard connectivity theorem of Bj\"orner, Lov\'asz, Vre\'cica,
and \v{Z}ivaljevi\'c \cite{BLVZ94} is proved, with its homological
prerequisite, in Appendix~\ref{sec:topology}:
\begin{equation}\label{eq:chess-vanish}
 N\ge2s-1
 \quad\Longrightarrow\quad
 \widetilde H_{s-2}(\Delta_{s,N};\F)=0.
\end{equation}

\begin{theorem}[Extension of prescribed matching moments]\label{thm:moments}
Suppose \(m\ge B\), \(N\ge2B-1\), and \(0\le k\le B\).
Every functional annihilating
\(\Ideal_k(\mathcal F^{\mathrm{fun}}_{m,N})\) extends to one annihilating
\(\Ideal_B(\mathcal F^{\mathrm{fun}}_{m,N})\).
The constant moment is preserved. In particular normalized functionals
exist through degree \(B\).
\end{theorem}
\begin{proof}
Use Lemma~\ref{lem:marginals}. Starting from the prescribed moments, fill
degrees \(s=k+1,\ldots,B\).
At degree one choose moments in each row whose sum is the prescribed
constant moment; a column exists because \(N\ge1\).

For \(s\ge2\), fix a set \(S\) of \(s\) rows. Form the reduced
\((s-2)\)-chain
\[
 z_S=\sum_{\substack{T\text{ a matching on }S\times[N]\\|T|=s-1}}
             z_T[T]
\]
in \(\Delta_{s,N}\).
This is a cycle. At \(s=2\), its augmentation is the sum of two equal
row totals, which is zero in \(\F\).
For \(s>2\), a face \(U\) of size \(s-2\) misses two rows of \(S\).
Its boundary coefficient is the sum of the prescribed marginal sums for
those two rows, hence \(z_U+z_U=0\).

Since \(N\ge2s-1\), equation~\eqref{eq:chess-vanish} gives an
\((s-1)\)-chain \(y_S\) with boundary \(z_S\).
Assign the coefficients of \(y_S\) as the new moments on full
\(s\)-matchings of \(S\times[N]\).
For a fixed \((s-1)\)-matching, its boundary equation is exactly the
missing-row equation in \eqref{eq:marginals}.

Different \(s\)-row sets have disjoint new unknowns: a size-\(s\) matching
has a unique occupied row set. Their shared lower faces retain their
already assigned values. Thus these fillings are globally consistent and
complete the induction.
The argument works for either value of the constant moment.
Starting with \(z_\varnothing=1\) at degree zero supplies a normalized
functional.
\end{proof}

\begin{corollary}[Old filtration stability and PC closure]\label{cor:stability}
Under \(m\ge B\) and \(N\ge2B-1\),
\begin{equation}\label{eq:old-stability}
 \Ideal_B(\mathcal F^{\mathrm{fun}}_{m,N})\cap\Pol_{\le k}
       =\Ideal_k(\mathcal F^{\mathrm{fun}}_{m,N})
       \quad(0\le k\le B),
 \qquad
 \Cons_B(\mathcal F^{\mathrm{fun}}_{m,N})
       =\Ideal_B(\mathcal F^{\mathrm{fun}}_{m,N}).
\end{equation}
In particular \(1\notin\Cons_B(\mathcal F^{\mathrm{fun}}_{m,N})\).
\end{corollary}
\begin{proof}
If a polynomial of degree at most \(k\) is outside \(\Ideal_k\),
finite-dimensional linear duality gives a functional vanishing on
\(\Ideal_k\) but not on that polynomial. Extend it by
Theorem~\ref{thm:moments}. The polynomial is then outside \(\Ideal_B\)
as well. The reverse inclusion in the first equality is immediate.

To prove PC closure, induct over a degree-\(B\) derivation.
Axioms and linear combinations lie in \(\Ideal_B\).
At a nonzero multiplication step \(g\mapsto X_{ij}g\), the predecessor
has ordinary degree at most \(B-1\). If \(g\in\Ideal_B\), the first
equality puts it in \(\Ideal_{B-1}\). Multiplication of its legal
generator multiples by \(X_{ij}\) stays in \(\Ideal_B\).
Hence \(\Cons_B\subseteq\Ideal_B\); the other direction is
Lemma~\ref{lem:pc-product}.
A normalized annihilator excludes one.
\end{proof}

The equalities in \eqref{eq:old-stability} concern the old functional
system. We will not use an analogous assertion for a system augmented
with extension blocks.
